\documentclass[11pt,a4paper]{article}
\usepackage[utf8]{inputenc}
\usepackage[T1]{fontenc}
\usepackage{amsmath, amssymb, amsthm}
\usepackage{geometry}
\usepackage{listings}
\usepackage{xcolor}
\usepackage{hyperref}

\geometry{margin=1in}

\title{Mathematical Foundations of the Transmutation Engine: \\ A Study in Computational Morphisms}
\date{\today}

\begin{document}

\maketitle

\begin{abstract}
This document provides an exhaustive mathematical derivation of the Transmutation Engine pipeline. The engine facilitates a series of representational shifts, transforming musical notation into quaternion rotations, evolving these rotations via cellular automata, and finally projecting the resultant states as holographic interference patterns rendered in ASCII space. This synthesis leverages techniques from high-dimensional rotation systems, discrete dynamical evolution, and wave optics.
\end{abstract}

\section{Introduction}
The Transmutation Engine is a transducer designed to explore the structural integrity of information across disparate mathematical domains. The pipeline follows the sequence:
\[ \text{Music} \xrightarrow{\mathcal{T}_1} \mathbb{H} \xrightarrow{\mathcal{T}_2} \text{CA} \xrightarrow{\mathcal{T}_3} \text{Hologram} \xrightarrow{\mathcal{T}_4} \text{ASCII} \]

\section{Stage 1: Harmonic to Quaternion Mapping ($\mathcal{T}_1$)}
Musical input is defined as a triplet $(f, d, a)$ representing frequency (Hz), duration (s), and amplitude (normalized). To preserve the orientation of these gestures in 4D space, we map them to a quaternion $q \in \mathbb{H}$.

\subsection{Normalization and Phase}
Given a reference frequency $f_{ref} = 440$ Hz, the normalized frequency is $n_f = f/f_{ref}$. The rotation angle $\theta$, representing the phase accumulated over the duration, is:
\[ \theta = 2\pi \cdot n_f \cdot d \]

\subsection{Rotation Geometry}
The rotation axis $\mathbf{v} = (v_x, v_y, v_z)$ is seeded by the harmonic structure and amplitude:
\[ v_x = \sin(n_f \pi), \quad v_y = \cos(n_f \pi), \quad v_z = a \]
To ensure $q$ lies on the unit hypersphere $S^3$, the axis is normalized:
\[ \hat{\mathbf{v}} = \frac{\mathbf{v}}{\|\mathbf{v}\|} \]

\subsection{Quaternion Construction}
Using the half-angle representation:
\[ q = \cos\left(\frac{\theta}{2}\right) + \hat{\mathbf{v}}\sin\left(\frac{\theta}{2}\right) \]
In scalar-vector form $q = (w, x, y, z)$:
\[ w = \cos\left(\frac{\theta}{2}\right), \quad x = \hat{v}_x \sin\left(\frac{\theta}{2}\right), \quad y = \hat{v}_y \sin\left(\frac{\theta}{2}\right), \quad z = \hat{v}_z \sin\left(\frac{\theta}{2}\right) \]

\section{Stage 2: Cellular Automata Evolution ($\mathcal{T}_2$)}
The CA serves as a discrete dynamical system that evolves the quaternion data through time.

\subsection{Seeding}
The initial state $C_0$ is a vector of cells of width $W=256$. Each quaternion $q_i$ occupies 4 spatial zones to preserve its full 4D dimensionality:
\[ c_{4i+k} = \lfloor 127.5(q_{i,k} + 1) \rfloor \quad \text{for } k \in \{w, x, y, z\} \]

\subsection{Rule 78 Dynamics}
Rule 78 is an elementary cellular automaton. In binary, 78 is $01001110_2$. For a cell $c_i$ at time $t$:
\[ c_i^{t+1} = \Phi(c_{i-1}^t, c_i^t, c_{i+1}^t) \]
Rule 78 was selected for its symmetric properties which maintain ordered structural fringes rather than chaotic entropy.

\section{Stage 3: Holographic Interference Projection ($\mathcal{T}_3$)}
The evolution of the CA over $T=64$ generations creates a space-time matrix. This is treated as the Object Wave ($O$) in a holographic simulation.

\subsection{The Interference Equation}
We simulate the intensity $I$ of a hologram resulting from the superposition of an Object Wave and a coherent Reference Wave ($R$):
\[ I(x, y) = |R(x, y) + O(x, y)|^2 \]

The Reference Wave is defined as:
\[ R(x, y) = \sin(k_x x + k_y y) \]
where $k_x=0.3$ and $k_y=0.1$ represent spatial frequencies. The Object Wave is derived from the normalized CA state:
\[ O(x, y) = \frac{c_x^y}{255} \]

Expanding the intensity:
\[ I(x, y) = R^2 + O^2 + 2RO \]
The cross-term $2RO$ encodes the interference fringes, producing the characteristic periodic visual morphology.

\section{Stage 4: ASCII Synthesis ($\mathcal{T}_4$)}
The intensity field $I \in [0, 1]$ is quantized against a character ramp $\Gamma$:
\[ \text{Char}(x, y) = \Gamma \left[ \lfloor I(x, y) \cdot (|\Gamma|-1) \rfloor \right] \]
where $\Gamma = \{ \text{' ', '.', ':', '-', '=', '+', '*', '\#', '\%', '@'} \}$.

\begin{thebibliography}{9}
\bibitem{uqgvs} 
Universal Quaternion Gesture Visual System (UQGVS), Technical Documentation.
\bibitem{fdca} 
FDCA-Simulator, Discrete Dynamical Systems Documentation.
\bibitem{holographic} 
Body-Holographic, Information Encoding Documentation.
\bibitem{wolfram} 
Stephen Wolfram, \textit{A New Kind of Science}, Wolfram Media, 2002.
\end{thebibliography}

\end{document}
